\documentclass[11pt]{article}
\usepackage[margin=0.9in]{geometry}
\usepackage{fontspec}
\usepackage{amsmath,amssymb,amsthm}
\usepackage{xurl}
\begin{document}\sloppy
\begin{center}
{\Large GROK H1144: cube 69 at 204 of 210 subcubes UNSAT; the closure checks read the catchup guard's timeout lines correctly; predictions held}\\[4pt]
2026-10-06 11:44 America/Los\_Angeles \quad branch \texttt{main}, tip \texttt{5115c782} \quad automation 41774ea2\\[2pt]
\textbf{P versus NP: OPEN.}
\end{center}

\section*{Peer delta}
Checked 11:44\,PT and again before commit. One peer commit landed after H1135's \texttt{96ae87f0}: \texttt{5115c782} (11:47\,PT, ``Add Sigma2 Sigma3 and functional hierarchy comparison track''), which adds only \path{research/explorations/search-spine/2026-10-06_SIGMA2_SIGMA3_FUNCTIONAL_TRACK.md}; content unread. Nothing changed under \path{research/main/} or \path{research/agents/grok/} beyond H1135. Root \texttt{main.tex} is untouched.

\section*{1. Literature}
None. The next claim is about the search pipeline.

\section*{2. Claims recorded}
\textbf{Timeout-format check (COMPUTATIONAL, exact on synthetic lines).} The H2340 guard \path{phaseB_catchup.sh} logs a failed subcube as \texttt{cube=$k$ sub=$s$ TIMEOUT\_or\_exit=$e$}, a different word from the direct runner's \texttt{TIMEOUT}, and it writes \texttt{CUBE-DONE $k$} even when some subcubes failed. The new script \path{fmt_check.py} (SHA-256 \texttt{ea3bfbbc\ldots}; it reads no log) feeds three synthetic lines to the timeout regexes of \path{cov.py} and \path{closure_run.py}: the guard's timeout line (exit 124) and a crash line (exit 1) both match with the right cube and subcube, and a non-timeout line does not. It also confirms that \path{closure_run.py} needs 210 distinct UNSAT subcube indices and never reads \texttt{CUBE-DONE}. So a cube-69, 72 or 74 subcube that times out or crashes keeps its cube open and cannot be counted as closed.

\textbf{Cube 69 (COMPUTATIONAL, solver-trusted, partial).} At 11:50\,PT the split log held 204 distinct cube-69 subcube indices, all UNSAT, with no SAT line and no cube-69 timeout; \texttt{CUBE-DONE 69} had not appeared. Cubes 72 and 74 had 0 lines. \path{cov.py} (unchanged, \texttt{e7296187\ldots}) again reports the open set $\{69,72,74\}$, all queued. First-step totals are unchanged: 72 of 75 kept cubes closed, 3 open, 0 SAT.

\section*{3. Held pre-registrations}
\path{oos69.py} (\texttt{af96f29c\ldots}), \path{oos72.py} (\texttt{012e135e\ldots}) and \path{oos74.py} (\texttt{d5e8d7e2\ldots}) were re-hashed and match H1115. None was run. No cube-69, 72 or 74 timing was read, and which cube-69 subcubes are still missing was not looked at, since that would leak timing into the cube-69 test.

\section*{4. What this changes}
Nothing about the bounds. They stay $9\le L(\texttt{a40ff54e580f5d0a},2)\le 11$, with the upper bound COMPUTATIONAL and the lower bound CONDITIONAL on the tightness lemma as recorded earlier. $L=9$ stays open while 3 cubes are open.

\section*{Effect on the standing candidate}
No claim is falsified and no charged resource improves. The time-state-representation surface stays CANDIDATE. The SAT-to-charged-model compiler and the product-span obstruction are both OPEN.

\section*{Next smallest step}
When \texttt{CUBE-DONE 69} appears, run \path{closure_run.py}; only if it lists 69 as closed, re-hash and run \path{oos69.py} unchanged. Then the same for cubes 72 and 74, and rerun \path{cov.py} after \texttt{CATCHUP-DONE}.

\bigskip
\noindent\textbf{Status sentence.} P versus NP remains OPEN: the SAT-to-charged-model compiler and the product-span obstruction are open load-bearing arrows. This pass records cube 69 at 204 of 210 subcubes UNSAT with 0 SAT and 0 timeouts (COMPUTATIONAL, solver-trusted, partial), checks that the closure scripts treat the catchup guard's timeout and crash lines as open (COMPUTATIONAL, synthetic lines), keeps \texttt{oos69.py}, \texttt{oos72.py} and \texttt{oos74.py} held with hashes unchanged, and leaves $9\le L(f,2)\le 11$ unchanged. No claim is falsified, no charged resource improves, and this is not a proof of anything about P versus NP.
\end{document}
